\documentclass[10pt,a4paper]{amsart}
\usepackage[margin=19mm]{geometry}
\usepackage{amsmath,amssymb,mathtools}
\usepackage[scr=rsfs,cal=euler]{mathalfa}
\usepackage{xfrac}
\usepackage[unicode,hidelinks,bookmarks=false]{hyperref}
\newcommand{\ie}{i.e.\ }
\newcommand{\cf}{cf.\ }
\newcommand{\resp}{respectively\ }
\newcommand{\Subsection}{\subsection}
\providecommand{\nobreakdash}{\nobreak\hbox{-}}
\setlength{\emergencystretch}{2em}
\multlinegap0pt
\usepackage{bm}
\usepackage{bbm}
\usepackage{dsfont}
\usepackage{xspace}
\usepackage{cancel}
\usepackage{mathtools}
\usepackage{amsthm}
\makeatletter
\@ifundefined{theorem}{\newtheorem{theorem}{Theorem}}{}
\@ifundefined{prop}{\newtheorem{prop}[theorem]{Proposition}}{}
\makeatother
\makeatletter
\newcommand{\R}{\mathbb{R}}
\expandafter\ifx\csname tagproblem\endcsname\relax
\newenvironment{tagproblem}[1]
 {\renewcommand\thetagprox{#1}\tagprox}
 {\endtagprox}
\fi
\makeatother
\title[Heteroclinic connections for fractional Allen-Cahn equations]{Heteroclinic connections for fractional Allen-Cahn equations with degenerate potentials}
\author{Francesco De Pas, Serena Dipierro, Mirco Piccinini, Enrico Valdinoci}
\date{Source: arXiv:2505.20054v1 (CC BY 4.0)}
\begin{document}
\maketitle

\noindent\emph{Source and licence.} Heteroclinic connections for fractional Allen-Cahn equations with degenerate potentials. Version arXiv:2505.20054v1 (CC BY 4.0), distributed under the Creative Commons Attribution 4.0 International licence, as recorded in the arXiv licence metadata for this exact version and shown on the abstract page https://arxiv.org/abs/2505.20054v1. Full rights notice: licenses/arxiv-2505.20054-CC-BY-4.0.txt.\par\smallskip

\noindent\textit{Editorial scope.} Counted unit: the Prop whose source label is \texttt{tonto} in the article (source0.tex lines 1071--1091, source label \texttt{tonto}), delivered together with the article's own complete original proof of it (source0.tex lines 1093--1235, one \texttt{proof} environment of 143 lines). No mathematics is written, repaired or re-derived: statement and proof are the article's own text line for line. Required context retained uncounted: krn\_symm (lines 211--213); main\_ellipt (lines 215--216). Every definition, display and label of those lines is retained unchanged. Omitted from this package and not redistributed: 1--210, 214--214, 217--1070, 1092--1092, 1236--3496. Nothing inside a retained excerpt is omitted or altered: every retained body is delivered exactly as the article's lines are written, so no omission receipt of any kind accompanies this package. Citations: the retained excerpts carry no citation command at all, so no bibliography entry of the article is transcribed and no reference is redistributed. Reference adaptation: none was needed and none was made; every reference inside the delivered unit resolves inside it, so no unresolved reference or citation is produced. Printed slips kept literal: no printed word, symbol, hypothesis or number of the counted statement or of its proof was altered. Layout and class: the article's own class is \texttt{amsart} with packages including inputenc, verbatim, lmodern, textcomp, amssymb, amscd, amsthm, amsrefs, amsxtra, mathtools, amsmath, dsfont, latexsym, mathrsfs; the wrapper is the lane's pinned \texttt{amsart} editorial wrapper at 10pt on A4, which restates the theorem-like environments the retained text needs and adds the article's own macro definitions that the retained text needs (\texttt{\textbackslash R}), with extra wrapper packages (bm, bbm, dsfont, xspace, cancel, mathtools). The delivered numbering is the unit's own. Assets: no retained span contains a figure environment, a graphics-inclusion command, a PDF-page inclusion, a table or a TikZ picture; no image file is copied into this package, the package declares no asset dependency and no image file is redistributed with it. Licence: version arXiv:2505.20054v1 of this article is distributed under the Creative Commons Attribution 4.0 International licence, as recorded in the arXiv licence metadata for this exact version and shown on the abstract page https://arxiv.org/abs/2505.20054v1. A full rights notice is delivered at licenses/arxiv-2505.20054-CC-BY-4.0.txt. Characters: every retained excerpt is pure ASCII, so the delivered engine, XeLaTeX with its default Latin Modern text font, reports no missing character.
\begin{align}
K(x)= K(-x) \qquad \text{for a.~\!e.}~x\in \R^n\,, \label{krn_symm}\tag{K1}
\end{align}

\begin{align}
\frac{\lambda\mathds{1}_{B_{r_0}}(x)}{|x|^{n+2s}}\leq K(x) \leq \frac{\Lambda}{|x|^{n+2s}} \qquad \text{for a.~\!e.}~x \in \R^n, \label{main_ellipt}\tag{K2} \end{align}

\begin{prop}\label{tonto}
Let~$\kappa\in (0, +\infty)$ and~$\sigma$, $\tau\in(1,+\infty)$. Let~$K$ satisfy~\eqref{krn_symm} and~\eqref{main_ellipt}. 

Let~$\phi \in C^{\infty}(\R)$ be such that
\begin{equation*}
\phi(x)= \begin{cases}
|x|^{-\sigma} &\mbox{if } x<-\kappa, \\
|x|^{-\tau} &\mbox{if } x>\kappa
\end{cases}
\end{equation*}
and
\begin{equation}\label{addipotesi76}
{\mbox{$\phi(x)\ge\gamma$ for all~$x\in[-\kappa,\kappa]$, for some~$\gamma\in(0,+\infty)$.}}
\end{equation}

Then,
\begin{equation}\label{yuuy_secs}
\lim_{|x | \to  +\infty} |x|^{1+2s} {L}_K\phi(x) \leq \Lambda \left( \frac{\kappa^{1-\sigma}}{\sigma-1}+ \int_{-\kappa}^{\kappa} \phi(y) \, dy+\frac{\kappa^{1-\tau}}{\tau-1}\right),
\end{equation}
where~$\Lambda$ is the quantity appearing in~\eqref{main_ellipt}.
\end{prop}

\begin{proof}
We will only prove~\eqref{yuuy_secs} as~$x \to -\infty$, being
the limit as~$x\to +\infty$ analogous.

For any~$x<-2\kappa$, we set
\begin{equation}\label{52BIS}
\begin{split}
&A_K(x):= PV_x \int_{-\infty}^{-\kappa}\left( \phi(y)-\phi(x)\right) K(x-y)\, dy, \\ &B_K(x):= \int_{-\kappa}^{\kappa} \left( \phi(y) -\phi(x)\right) K(x-y) \, dy\\
{\mbox{and }}\qquad
&D_K(x):= \int_{\kappa}^{+\infty}\left( \phi(y)-\phi(x)\right) K(x-y)\, dy.
\end{split}
\end{equation}
In the following computations we will omit the principal value notation,
for the sake of readability.

We claim that 
\begin{equation}\label{x_secs}
\lim_{x \to -\infty} |x |^{1+2s} A_K(x) \leq  \frac{\Lambda\kappa^{1-\sigma}}{\sigma-1}.
\end{equation}
In order to show this, we apply the change of variable~$y:=|x|\theta$ and compute
\begin{equation*}
A_K(x) = \int_{-\infty}^{-\kappa} \left( | y |^{-\sigma} - |x |^{-\sigma} \right) K(x-y)\, dy 
= |x |^{1-\sigma} 
\int_{-\infty}^{-{\kappa}/{|x|}} \left( |\theta |^{-\sigma}-1\right) K\left( x (1+\theta)\right) \, d\theta.
\end{equation*}
We also set
\begin{equation}\label{53BIS}
\begin{split}
& A_{K,I}(x):= |x |^{1-\sigma} \int_{-\infty}^{-{3}/{2}}\left( |\theta |^{-\sigma}-1\right) K\left( x (1+\theta)\right) \, d\theta, \\
& A_{K,II}(x):=  |x |^{1-\sigma} 
\int_{-{3}/{2}}^{-{1}/{2}}\left(|\theta |^{-\sigma}-1\right) K\left( x (1+\theta)\right) \, d\theta \\ \mbox{and}\qquad
&A_{K,III}(x):=  | x |^{1-\sigma} \int_{-{1}/{2}}^{-{\kappa}/{| x|}}\left( | \theta |^{-\sigma}-1\right) K\left( x (1+\theta)\right) \, d\theta.
\end{split}
\end{equation}
Exploiting~\eqref{main_ellipt}, we obtain that
\begin{equation}\label{equno}
|A_{K,I}(x)| \leq \Lambda |x |^{-(\sigma+2s)}\int_{-\infty}^{-{3}/{2}}\frac{1- | \theta |^{-\sigma}}{| 1+\theta|^{1+2s}} \, d\theta \leq C |x |^{-(\sigma+2s)},
\end{equation}
for some~$C>0$, depending on~$s$, $\sigma$ and~$\Lambda$.

Moreover, we check that
\begin{equation}\label{jgsp}
| A_{K,II}(x)| \leq C |x |^{-(\sigma+2s)}.
\end{equation}
To this aim, we use the change of variable~$z:= 1+\theta$ and
the symmetry of~$K$ in~\eqref{krn_symm}
to see that
\begin{equation*}
\begin{split}
&\int_{-{3}/{2}}^{-{1}/{2}} \left( |\theta|^{-\sigma}-1 \right) K \left( x(1+\theta)\right) \, d\theta 
= \int_{-{1}/{2}}^{{1}/{2}} \left( ( 1-z)^{-\sigma}-1\right) K(xz) \, dz \\
&\qquad=\int_{-{1}/{2}}^{{1}/{2}} \left( \sigma z +O(|z |^2) \right) K(xz) \, dz 
=\int_{-{1}/{2}}^{{1}/{2}} O(|z |^2) K(xz) \, dz.
\end{split}
\end{equation*}
As a consequence,
\begin{equation*}
\begin{split}
&\left| \int_{-{3}/{2}}^{-{1}/{2}} \left( |\theta|^{-\sigma}-1 \right) K \left( x(1+\theta)\right) \, d\theta \right|
= \left|\int_{-{1}/{2}}^{{1}/{2}} O(|z |^2)  K(xz) \, dz\right|\\
&\qquad\qquad\leq C\Lambda |x |^{-(1+2s)}  \int_{-{1}/{2}}^{{1}/{2}}| z|^{1-2s}\, dz \leq C |x |^{-(1+2s)} ,
\end{split}
\end{equation*} up to renaming~$C>0$.
This leads to~\eqref{jgsp}, as desired.

Furthermore, observing that~$\theta^{-\sigma}>1$ if~$\theta \in (0,1/2)$, we obtain that
\begin{equation*}
\begin{split}
A_{K,III}(x) & \leq \Lambda |x |^{-(\sigma+2s)} 
\int_{-{1}/{2}}^{-{\kappa}/{|x|}} \frac{ |\theta|^{-\sigma}-1}{ (1+\theta)^{1+2s}} \, d\theta\\
&= \Lambda |x |^{-(\sigma+2s)} 
\left(
\int_{{\kappa}/{| x |}}^{{1}/{2}} 
\theta^{-\sigma}\big(1+O(\theta)\big) \, d\theta -
\int_{{\kappa}/{|x |}}^{{1}/{2}} \frac{d\theta}{(1+\theta)^{1+2s}}\right)\\
&= \Lambda |x |^{-(\sigma+2s)} 
\left( \frac{1}{1-\sigma}\left(2^{\sigma-1} -\kappa^{1-\sigma}|x|^{\sigma-1} \right) +O(|x|^{\sigma-2})
-\int_{{\kappa}/{|x |}}^{{1}/{2}} \frac{d\theta}{(1+\theta)^{1+2s}}\right).
\end{split}
\end{equation*}
Combining this with~\eqref{equno} and~\eqref{jgsp},
and noticing that~$A_K(x) = A_{K,I}(x)+ A_{K,II}(x)+ A_{K,III}(x)$,
we obtain~\eqref{x_secs}. 

Now, we want to show that
\begin{equation}\label{fcpmd}
\lim_{x \to - \infty} |x |^{1+2s} B_K(x) \leq  \Lambda \int_{-\kappa}^{\kappa} \phi(y)\, dy.
\end{equation}
For this, we take~$\gamma$ as in~\eqref{addipotesi76}
and we suppose that~$|x|>\gamma^{-1/\sigma}$.
In this way, we have that~$\phi(x)=|x|^{-\sigma}<\gamma\le\phi(y)$
for all~$y\in(-\kappa,\kappa)$,
thanks to~\eqref{addipotesi76}.
Therefore, by~\eqref{main_ellipt},
\begin{eqnarray*}
|x|^{1+2s}B_K(x)&=& |x|^{1+2s}
\int_{-\kappa}^{\kappa} \left( \phi(y) -\phi(x)\right) K(x-y) \, dy\\
&\le& \Lambda \int_{-\kappa}^{\kappa} \left( \phi(y)-\phi(x) \right) \left| 1 -\frac{y}{x} \right|^{-(1+2s)}\, dy.
\end{eqnarray*}
We can now take the limit as~$x\to-\infty$ and use the
Dominated Convergence Theorem to deduce~\eqref{fcpmd}.

Furthermore, we claim that
\begin{equation}\label{fcpmt}
\lim_{x \to - \infty} |x |^{1+2s} D_K(x) \leq  \frac{\Lambda\kappa^{1-\tau}}{\tau-1},
\end{equation}
Indeed, the change of variable~$y:=|x|\theta$ 
and~\eqref{main_ellipt} give that
\begin{equation*}
\begin{split}
D_K(x)&= \int_{\kappa}^{+\infty} (y^{-\tau} - |x |^{-\sigma})K(x-y) \, dy \\
&= |x |\int_{{\kappa}/{|x|}}^{+\infty}\big(
|x |^{-\tau} \theta^{-\tau} - |x |^{-\sigma}\big) K(x (1+\theta))\, d\theta\\
&\leq \Lambda |x |^{-2s} \int_{{\kappa}/{|x|}}^{+\infty} \frac{|x |^{-\tau} \theta^{-\tau} +|x |^{-\sigma}}{(1+\theta)^{1+2s}}\, d\theta\\
&\leq  \Lambda |x |^{-(\tau+2s)} 
\int_{{\kappa}/{|x|}}^{+\infty}  \frac{\theta^{-\tau}}{(1+\theta)^{1+2s}}\, d\theta
+\Lambda | x |^{-(\sigma+2s)}
\int_{{\kappa}/{| x|}}^{+\infty} \frac{d\theta}{(1+\theta)^{1+2s}}.
\end{split}\end{equation*}
We notice that
\begin{eqnarray*}
\int_{{\kappa}/{|x|}}^{+\infty}  \frac{\theta^{-\tau}}{(1+\theta)^{1+2s}}\, d\theta
&=&\int_{{\kappa}/{|x|}}^{1/2}
\theta^{-\tau}\big(1+O(\theta)\big)\, d\theta
+ \int_{1/2}^{+\infty}  \frac{\theta^{-\tau}}{(1+\theta)^{1+2s}}\, d\theta\\
&=&\frac{1}{1-\tau}\left( 2^{\tau-1}-\kappa^{1-\tau}|x|^{\tau-1}
\right) + O(|x|^{\tau-2})
+ \int_{1/2}^{+\infty}  \frac{\theta^{-\tau}}{(1+\theta)^{1+2s}}\, d\theta.
\end{eqnarray*}
As a result,
\begin{eqnarray*}
|x|^{1+2s}D_K(x)&\le& 
\Lambda|x|^{1-\tau}\left(\frac{1}{1-\tau}\left( 2^{\tau-1}-\kappa^{1-\tau}|x|^{\tau-1}
\right) + O(|x|^{\tau-2})
+ \int_{1/2}^{+\infty}  \frac{\theta^{-\tau}}{(1+\theta)^{1+2s}}\, d\theta\right)\\&&\qquad
+\Lambda |x |^{1-\sigma}
\int_{{\kappa}/{|x|}}^{+\infty} \frac{d\theta}{(1+\theta)^{1+2s}},
\end{eqnarray*}
which gives the desired claim in~\eqref{fcpmt}.

Thus, since~$L_K\phi(x)= A_K(x) + B_K(x) + D_K(x)$,
combining together~\eqref{x_secs}, \eqref{fcpmd} and~\eqref{fcpmt} yields the thesis.
\end{proof}    

\end{document}
